
\documentclass[11pt]{article}


\usepackage{graphicx}
\usepackage{enumitem}
\usepackage{placeins}
\usepackage{natbib}
\usepackage{booktabs}
\usepackage{amsmath,amssymb}


\usepackage[letterpaper,margin=1in]{geometry}
\usepackage{longtable}
\usepackage{array}
\usepackage{tabularx}
\usepackage{microtype}
\usepackage[hidelinks]{hyperref}
\usepackage{xurl}

\newcolumntype{Y}{>{\raggedright\arraybackslash}X}

% Make nested bullet lists less indented and more compact.
\setlist[itemize,1]{leftmargin=1.35em,label=\textbullet}
\setlist[itemize,2]{leftmargin=1.35em,label=--}
\setlist[itemize,3]{leftmargin=1.35em,label=*}
\setlist[itemize]{itemsep=0.35em,topsep=0.35em,parsep=0pt,partopsep=0pt}

% Helps prevent long model/dataset names from running off the page.
\setlength{\emergencystretch}{3em}
\newcommand{\model}[1]{\texttt{#1}}

\title{The Judge Is Not Impartial: Self-Preference in Medical LLM Evaluation}

\author{%
\begin{tabular}{c}
\textbf{Zara Ansari}\textsuperscript{1,2}
\hspace{1em}
\textbf{Shlok Natarajan}\textsuperscript{1}
\hspace{1em}
\textbf{Author Three}\textsuperscript{1}
\\[0.25em]
\textbf{Author Four}\textsuperscript{1}
\hspace{1em}
\textbf{Aaron Fanous}\textsuperscript{1}
\hspace{1em}
\textbf{Roxana Daneshjou}\textsuperscript{1,3}
\\[0.7em]
\small
\textsuperscript{1}Department of Biomedical Data Science,
\textsuperscript{3}Department of Dermatology
\\
\small
Stanford University, 450 Jane Stanford Way,
Stanford, CA 94305, USA.
\\
\small
\textsuperscript{2}Department of Computer Science,
Harvey Mudd College
\\
\small
301 Platt Boulevard, Claremont, CA 91711, USA.
\\[0.6em]
\small
Emails: zansari6@stanford.edu,
insert-email,
author2@stanford.edu
\\
\small
author3@stanford.edu,
author4@stanford.edu,
roxanad@stanford.edu
\end{tabular}%
}

\date{August 2026}

\begin{document}

\maketitle

\begin{abstract}

Large language models (LLMs) are increasingly used to perform medical tasks, creating the need for medical benchmarks to assess which model performs best. However, the scale of benchmarking makes human review untenable, and large language models (LLMs) are increasingly used to evaluate model-generated medical content, creating a potential circularity: a judge may show self-preference --- favoring outputs from the same model independently of clinical quality. To characterize LLM judge self-preference in medical settings, we measured same-model affinity in single-turn specialist question answering and
multi-turn standardized-patient encounters. We evaluated eight models, two from OpenAI, Anthropic, Google, and Qwen, on 620 Real-POCQi questions. The same eight models also generated and judged 200 MedSP1000 scenarios after two, four, six, and eight visible turns. Model identities were concealed and response order was randomized.

 While holding a model responses to medical queries fixed, we assessed how different LLM judges assessed the response and the degree of self-preference. In Real-POCQi, judges ranked their own answers 1.22 positions closer to first on average (95\% CI $[1.18,1.26]$). The self preference effect persisted after adjusting for presentation
position and differences in how judges used the score scale. Interestingly, self-preference increased when models were not provided a a grading rubric to follow. An exploratory family-level comparison found a similar pattern, although it
could not separate exact-model preference from broader family preference.
Revealing model names
modestly weakened rather than amplified the effect. In MedSP1000, own-model
ranks were better by 0.57, 0.53, 0.62, and 0.61 positions at two, four,
six, and eight turns, respectively; there was no clear overall difference
across conversation lengths ($p=0.17$). We found that Answer Rank Bias, Answer Score Bias, and Own-Answer Win Rate capture different aspects of self-preference
\iffalse
Answer Score Bias and Own-Answer Win Rate sometimes led to
different conclusions because the latter also reflected answer quality.
\fi
These findings support increased vigilance including fixed-answer analyses, multiple evaluation measures,
and judge panels spanning model families when deploying LLM-as-a-judge systems in medical contexts.

\end{abstract}

\section{Introduction}

As large language models (LLMs) are considered for different tasks in medicine, benchmarking has become a key method for choosing the best model for a particular task.  However, the complexity of medical practice has called for larger and more complex benchmarking to compare the performance of AI models. Human review of thousands of scenarios tested across multiple AI models is untenable due to cost and time \citep{zheng2023judging}. A recent review of medical AI benchmarking found that large language models (LLMs) are increasingly used as judges to evaluate the outputs of other LLMs (and in some cases, their own outputs) \citep{li2026llm}. This same study found that OpenAI's models are currently the most popular judges for medical AI benchmarking studies \citep{li2026llm}.

Researchers and industry experts claim using an LLM as a judge provides a scalable way to evaluate responses based on certain criteria such as accuracy, clinical utility, source quality, verifiability, and completeness.  The stakes in medical AI evaluation are high, as biased judges could incorrectly favor unsafe models and influence clinical deployment decisions to the detriment of patients. Furthermore, evaluating medical responses is especially complicated since an answer may sound confidently correct, but still contain incomplete medical information. As a result, evaluating LLMs for medical tasks is a multifaceted process where several dimensions of response quality must be considered simultaneously.

Prior work in general generative AI tasks has shown that LLM judges may exhibit biases such as positional and verbosity bias. \cite{wataoka2025selfpreference} examined how LLM judges may favor responses that reflect policies or writing styles familiar to them.  In fact, one study showed that LLM judges may favor AI generated clinical responses over human produced ones. \cite{alvarezarenas2026unreliable} found that this preference seemed to depend more on how an answer was written, especially how long it was and the language used, than on whether the model could actually tell if the answer came from a human or an AI. In one experiment, long responses still received high scores even when matched with unrelated questions, demonstrating that verbosity can outweigh the content for some LLM judges. \cite{wataoka2025selfpreference} introduced a quantitative way to measure self-preference in LLM judges. They found that GPT-4 showed significant self-preference and that LLM judges tended to rate lower perplexity responses more highly than human evaluators did.  Work from \cite{pombal2026selfpreference} looked at whether rubric-based evaluations (with a set of yes or no criteria) could mitigate the self-preference of LLM judges. Their work found that LLM self-preference persists even with fully "objective" rubrics.

The literature for self-preference in medical contexts is far more sparse. To address these gaps, we evaluate self-preference across four model families in single-turn generalist-specialist question answering and in the multi-turn patient–physician conversations. In order to address the fact that a model may choose its own response because its answer is genuinely better, we evaluate the models preference in three distinct ways. Our primary measure is Answer Rank Bias, which compares how a model's own judge and outside judges rank the same fixed response. We also use Answer Score Bias, which quantifies how different judges score the same fixed response. Furthermore, we use Own-Answer Win Rate, which measures how often a judge ranks its own model's response above a competing response. Our findings of LLM judge self-preference in medical contexts provide important considerations for future medical AI benchmarking studies.


%\section{Related Work}

%\subsection{AI-Generation Preference in LLM Clinical Judges}

%Recent work has shown that LLM judges may favor AI generated clinical responses over human produced ones. \cite{alvarezarenas2026unreliable} found that this preference seemed to depend more on how an answer was written, especially how long it was and the language used, than on whether the model could actually tell if the answer came from a human or an AI. In one experiment, long responses still received high scores even when matched with unrelated questions. This shows that verbosity can outweigh the relevance of content for some LLM judges. Their activation steering experiments showed that this preference disappeared when the model was adjusted to rely less on response length when judging answer quality. However, this work compares human and AI generated responses rather than testing whether a judge favors outputs from its own model family. Our study examines this question directly in pairwise medical response evaluation, testing whether models give an unfair advantage to answers generated by their own family compared with other AI models.


%\subsection{Self-Preference in Rubric-Based Evaluation}
%\cite{pombal2026selfpreference} study whether LLM judges evaluate answers from their own model family more favorably. Their main approach uses rubric-based evaluation, where a judge looks at one answer and decides whether it meets a set of specific yes or no criteria. They also test this on HealthBench, which is a medical benchmark that uses clinical rubrics. They evaluate 12 models spanning five model families as judges and generators. On HealthBench, they found that self-preference bias can shift a model's score by up to 10 points. Because leading models often perform within only a few points of one another, bias of this magnitude can substantially alter their relative rankings. They also found that self-preference is more pronounced for negative rubrics. These rubrics identify harmful or undesirable behaviors that a response should avoid, such as making a clinical error or omitting an important warning. Self-preference was also stronger for more subjective topics, such as emergency referrals and patient communication. Although they compare rubric-based evaluation with other formats, including pairwise comparisons, their main focus is on evaluating static responses. Our study focuses on pairwise medical evaluation and also includes multi-turn physician-patient conversations. In these conversations, each physician model can ask different follow up questions and guide the interaction in different ways. This can change what the patient shares later. This lets us test whether self-preference affects how judges evaluate the quality and safety of a full clinical interaction.


%\subsection{Self-Preference in Pairwise LLM Evaluation}
%Prior work has shown that LLMs may favor responses generated by their own model family during pairwise evaluation.
%\cite{panickssery2024llm} examine this relationship using news summarization tasks from the XSUM and CNN/DailyMail datasets. In their study, models compare summaries of the same article written by humans or different LLMs to test whether they can recognize their own writing and whether that recognition leads to self-preference. Our study examines the same issue in the medical domain. Instead of comparing which summary best represents an article, the judges compare two physician-style responses to a patient question based on factors such as safety, completeness, clarity, and medical appropriateness.


%\subsection{Self Preference Bias in LLM Judges}
%\cite{wataoka2025selfpreference} introduced a quantitative way to measure self-preference in LLM judges. They found that GPT-4 showed significant self-preference and that LLM judges tended to rate lower perplexity responses more highly than human evaluators did. This suggests that models may prefer responses that feel more familiar them. However, their study focused on general domain dialogue. Our study extends this work to the medical domain by testing four model families on both single-turn medical answers and multi-turn patient–physician conversations. We also examine self-preference at both the score level and the decision level, which allows us to compare how generously models score their own responses with how often they choose them in direct comparisons.












\section{Methods}


%\subsection{Dataset Description}

Real-POCQi contains 620 real world clinical questions which were submitted by practicing U.S. physicians through the OpenEvidence platform. The questions span 30 medical specialties and focus on point-of-care clinical decision making. Each entry includes a unique question ID, the question text, and its respective specialty. The dataset was reviewed to remove identifying patient information and other content that could raise privacy concerns. In this study, we use all 620 questions as prompts for generating and evaluating responses from the eight tested models.

MedSP1000 contains standardized-patient scenarios derived from published clinical teaching cases. Each scenario includes separate materials for the physician and the standardized patient. The dataset also includes evaluator and environment-controller materials, which we exclude because our study focuses on generating and evaluating the consultations. We screened the scenarios for standard patient–physician format, sufficient patient information, English-language content, and the absence of diagnosis-key leakage, caregiver-proxy framing, or preverbal patients, leaving 217 eligible scenarios. From here, we randomly sampled 200 scenarios for our analysis. In this study, we provide the physician-facing information to the simulated clinician model and the patient information to guide the simulated patient model, generating multi-turn patient–physician conversations across the eight tested models. The patient model was held as Mistral-Small-3.1-24B-Instruct-2503 across all runs to avoid any potential same model family bias.

\subsection{Sample Size Analysis}

We conducted a power analysis using a two-sided significance level of 0.05 and 90\% power. For the Answer Score Bias analysis, described in the following sections, approximately 119 paired observations were needed to detect a small-to-medium effect ($d=0.30$). For the Own-Answer Win Rate analysis, approximately 114 observations were needed to detect a win rate of 65\%, compared with the 50\% rate expected under no preference. Based on these estimates, the 620 questions used in the single-turn setting and the 200 scenarios used in the multi-turn setting provided more than enough observations for both analyses.

In the single-turn setting, each of the eight models answered all 620 questions, resulting in 4,960 specialist responses (620 questions $\times$ 8 models). These responses were then evaluated listwise. All eight models also served as judges, with each judge scoring and ranking the full set of eight candidate responses for every question. This resulted in 4,960 listwise judgments (8 judges $\times$ 620 questions). In each judgment, the judge scored all eight responses and then ranked them, with the ranking not required to follow the rubric scores.

In the multi-turn setting, each of the eight models generated an eight-turn conversation for all 200 scenarios, producing 1,600 conversations (200 scenarios $\times$ 8 models). The shorter 2, 4, and 6 turn versions were created by truncating these eight-turn conversations. Although, regenerating these conversations entirely for separate turn counts coudl be an area of further study. Listwise judging was then performed separately at each of the four conversation lengths using the same eight judges, resulting in 6,400 judgments (8 judges $\times$ 4 lengths $\times$ 200 scenarios). In total, our analysis includes 11,360 identity-blind, order-randomized listwise judgments across the two settings. The no-rubric and identity-revealed checks add 800 and 1,200 judgments, bringing the study total to 13,360.

\subsection{Generator and Judge Models}
The study uses two models from each of four families, shown in
Table~\ref{tab:cohort}. Generator models are fed the original query from either the Real-POCQi or MedSP1000 dataset and produce an answer in response to that given query. In the single turn and multi turn pipelines, all eight models act as generators and judges.


These tier labels describe their role in this study and
are not intended as a universal performance ranking.

\begin{table}[htbp]
\centering
\caption{Generator and judge cohort.}
\label{tab:cohort}
\begin{tabular}{lll}
\toprule
Family & Higher-capability model & Workhorse model \\
\midrule
OpenAI & \model{gpt-5.6-sol} & \model{gpt-5.6-terra} \\
Anthropic & \model{claude-opus-5} & \model{claude-sonnet-5} \\
Google & \model{gemini-3.1-pro-preview} & \model{gemini-3.7-flash} \\
Qwen & \model{Qwen/Qwen3.5-122B-A10B-FP8} & \model{Qwen/Qwen3.8-27B-FP8} \\
\bottomrule
\end{tabular}
\end{table}
\FloatBarrier

\subsection{Single-turn Pipeline}

This pipeline evaluates specialist model behavior on single-turn questions using all 620 questions from the Real-POCQi dataset, frozen at source revision 9002e1dd. Each question is associated with a clinical specialty, such as cardiology or neurology. For each question, eight models (GPT-5.6 Sol, GPT-5.6 Terra, Claude Opus 5, Claude Sonnet 5, Gemini 3.1 Pro, Gemini 3.7 Flash, Qwen3.5-122B, and Qwen3.8-27B) each independently generate a single answer in the role of the relevant specialist, prompted with that specialty's label. The eight models represent two capability tiers from each of four model families. Each model generates exactly one response per question, resulting in 4,960 total answers.

These responses are evaluated in a listwise manner across five criteria: accuracy, clinical utility, source quality, verifiability, and completeness (0–5; overall = their mean). For each question, a judge receives all eight responses at the same time. The responses are identity-blind, and their presentation order is randomized independently for each judge and question under a fixed seed. In the same evaluation, the judge also ranks all eight responses from best to worst using its overall clinical judgment. The ranking is not required to follow the rubric scores, and rankings are strict with no ties. All eight models also serve as judges, resulting in 4,960 judgments. We refer to this as the primary condition, in which all 620 questions are judged with both rubric scores and a full ranking, and generator identities are hidden. These judgments yield three measures of self-preference. Answer Rank Bias is our primary measure, and we also report Answer Score Bias and Own-Answer Win Rate.

We also evaluated two sensitivity conditions, reusing the same generated responses so that only the judging setup changes. The no-rubric condition tests whether the five criteria has any effect. All eight judges rank the same blinded responses on overall quality alone, with no criteria provided, on a subset of 100 of the 620 questions. The identity revealed condition tests whether the effect depends on knowing who wrote each response (labels every candidate with its generator model) on a subset of 200 questions. Only six models serve as judges in the identity revealed condition, since the two self-hosted Qwen models were not run, but all eight generators' responses are still ranked.


\subsection{Multi-turn Pipeline}

This pipeline evaluates physician model behavior in multi-turn patient–physician conversations, using 200 scenarios from MedSP1000. Each scenario provides case material written for the examinee, which we give to the physician model. A separate private set of instructions written for the standardized patient, which we give to the patient model. Format and level of detail vary across cases. The physician only sees its own material and the patient's replies, and starts the consultation.

Eight physician models, GPT-5.6 Sol, GPT-5.6 Terra, Claude Opus 5, Claude Sonnet 5, Gemini 3.1 Pro, Gemini 3.7 Flash, Qwen3.5-122B, and Qwen3.8-27B, each independently generate an 8 turn conversation for every scenario. Physician turns are produced by the model under study, while patient turns are generated by a standardized patient simulator, Mistral Small 3.1 24B Instruct, which is not one of the evaluated models. Each physician model is given four speaking turns, resulting in four physician–patient exchanges. Each model generates exactly one conversation per scenario, producing 1,600 total conversations. Shorter conversation lengths of 2, 4, and 6 turns are derived by truncating each 8 turn conversation to the corresponding length. Since each turn depends only on the conversation history that comes before it, truncating the longer conversations provides valid shorter versions.

Similar to the single-turn pipeline, conversations in the multi-turn pipeline are evaluated in a listwise manner across the same five criteria (0–5; overall = their mean). These criteria are applied to the physician’s behavior across the encounter. For each scenario, a judge receives all eight conversations at the same time and scores only the physician responses, while the patient turns are included as context but are not scored. The conversations are identity-blind, and their presentation order is randomized independently for each judge and scenario under a fixed seed. In the same evaluation, the judge also ranks all eight conversations from best to worst using its overall clinical judgment. The ranking is not required to follow the rubric scores, and rankings are strict with no ties. All eight physician models also serve as judges.

The difference from the single-turn pipeline is that this listwise judging is performed separately at each of the four conversation lengths (2, 4, 6, and 8 turns). This results in 1,600 judgments at each conversation length and 6,400 judgments in total. Thus, the listwise judgments and the three resulting measures of self-preference are computed and reported separately for each model and conversation length.


\section{Analysis of Data}

We summarize self-preference using three complementary measures: \textit{(1) Answer Rank Bias (ARB)}, our primary measure, holds the exact response fixed and compares the rank assigned by its own-model judge with the average rank assigned by outside judges. Negative values indicate greater self-preference (models ranking their own responses in lower-numbered positions relative to others). \textit{(2) Answer Score Bias (ASB)} makes the analogous fixed-response comparison using mean rubric scores. Positive values indicate that the judge scores its response more generously. \textit{(3) Own-Answer Win Rate (OAWR)} is the proportion of implied head-to-head comparisons in which a judge ranks its own model's response above a competitor. Unlike the first two measures, Own-Answer Win Rate compares different responses and can therefore reflect genuine differences in response quality as well as judge preference.

\subsection{Answer Rank Bias}

Every judge ranks the full candidate list, so each answer receives $J$ ranks: one from the judge belonging to the model that wrote it, and $J-1$ from outside judges. This holds in every condition where all candidate models also serve as judges. For question $q$ and exact model $j$, let $r_{kqj}$ be the rank that judge $k$ assigns to the answer generated by model $j$, where lower ranks are better. Answer Rank Bias is
\begin{equation}
D_{jq}=r_{jqj}-\frac{1}{J-1}\sum_{k\ne j}r_{kqj}.
\end{equation}
The first term is the rank a model's own judge gives its own answer. The second is the average rank the outside judges give that same answer. So, $D_{jq}<0$ means that judge $j$ ranks its own answer more favorably than outside judges rank the same answer to the same question. Because every judge is ranking the same answer to the same question, the differences in answer quality cancel out. This means that a model that writes strong answers is not credited with self-preference. Conversely, a model whose answers rank poorly can still show self preference.

Where a candidate appears in the prompt can affect the rank it receives \citep{zheng2023judging}. We therefore estimate each position's effect, using judge-generator terms to keep those estimates unbiased. Only the position effect is removed before computing the difference. Each judge's value is computed partly from the other judges' ranks, so the $J$ values for a question are related rather than independent. We average them within each question and compute confidence intervals across questions. As a check we repeat the comparison on standardized scores, and where judges are tested separately $p$-values carry Holm correction.


The four MedSP1000 conditions show the same $1{,}600$ conversations at increasing length, so every conversation contributes a value at all four lengths. Length comparisons are thus paired within question. If we write $D^{(2)}$ through $D^{(8)}$ for the value of $D_{jq}$ at each length, then $D^{(4)}-D^{(2)}$ (4 minus 2 turns) is the change in self-preference for a question when four turns are visible relative to two. Testing the three changes relative to two turns separately gives three chances at a false positive. The joint length test asks in one step whether all three are zero on average, using Hotelling's $T^{2}$ test, which is the multivariate version of a paired $t$-test.


\subsection{Answer Score Bias}
Let $s_{kqj}$ represent the mean of judge $k$’s five rubric scores for the answer generated by model $j$ on question $q$ (which uses the same $0$--$5$ scale as the criteria). To measure Answer Score Bias, we hold the answer fixed and compare the score assigned by the model’s own judge with the average score assigned by the outside judges:

\begin{equation}
B_{jq}=s_{jqj}-\frac{1}{J-1}\sum_{k\ne j}s_{kqj}.
\end{equation}

Because all judges evaluate the same response, differences in the quality of the answer itself are held constant. The remaining difference reflects how the judges score that response. A positive value means the model’s own judge scores its answer more generously than the outside judges do. A negative value means it scores its own answer more harshly.

\subsection{Own-Answer Win Rate}

The listwise rankings contain no ties, so each ranking can also be treated as a set of head-to-head comparisons between the judge's own answer and every competing answer. Let $Q$ represent the number of questions in the condition. The Own-Answer Win Rate for model $j$ is
\begin{equation}
P_j=\frac{\sum_q\sum_{m\ne j}\mathbb{I}(r_{jqj}<r_{jqm})}{Q(J-1)},
\end{equation}
where the numerator counts the comparisons in which judge $j$ ranked its own answer above a competitor and the denominator is the total number of such comparisons, $Q(J-1)$. Unlike Answer Score Bias, this measure does not hold answer quality constant because the competing responses are different. So, a model may rank its own answer higher simply because that answer is genuinely better. The ranking is produced in the same evaluation as the rubric scores, but it is not required to
follow those scores.



\section{Results}

\subsection{Real-POCQi Generation}

All 4,960 intended answers were generated, so every answer was judged by all eight judges. Claude Opus 5 was much more verbose than the others, with a median of 2,927 output tokens against 1,354 for Claude Sonnet 5 and 656-988 for the remaining six models. A manual review of three questions, covering the four generators found factual errors (venous-reflux error), internal contradictions (HSV-trial contradiction), and unsupported claims across models (in an oncology question). Because our study measures how judges behave rather than how accurate the models are, this variation is what makes this usable. It gives the judges real quality differences to detect. It also means the answers should not be treated as a clinical gold standard.

Table~\ref{tab:real-performance} shows how the eight judges collectively ranked the generators. With eight candidates the average rank is 4.5, so values below 4.5 indicate above-average performance in the judges' collective rankings.

\begin{table}[htbp]
\centering
\small
\caption{Real-POCQi aggregate generator performance across eight
judges. Lower mean rank is better.}
\label{tab:real-performance}
\begin{tabular}{lrrr}
\toprule

Generator & Mean rank & Ranked first & Mean rubric sum \\
\midrule
Claude Opus 5 & \textbf{1.89} & \textbf{69.3\%} & \textbf{23.11} \\
GPT-5.6 Terra & 3.37 & 17.5\% & 21.65 \\
Gemini 3.7 Flash & 3.92 & 1.3\% & 21.08 \\
GPT-5.6 Sol & 4.30 & 8.3\% & 20.84 \\
Claude Sonnet 5 & 4.39 & 1.5\% & 20.63 \\
Gemini 3.1 Pro & 4.62 & 1.3\% & 20.41 \\
Qwen 3.5 122B & 6.52 & 0.4\% & 17.52 \\
Qwen 3.8 27B & 6.99 & 0.3\% & 16.41 \\


\bottomrule
\end{tabular}
\end{table}
\FloatBarrier

Claude Opus 5 ranked first by a wide margin on all three measures, and the two Qwen models ranked last. Mean rank and first-place rate can disagree. Gemini 3.7 Flash had a better mean rank than GPT-5.6 Sol but was ranked first far less often (1.3\% versus 8.3\%), so Flash is consistently good without often being best, while Sol is more uneven. The gap to the Qwen models is larger than any gap inside the top six, which span 2.7 rank positions, while a further 1.9 positions separate sixth place from seventh. Rubric scores make that gap look smaller. They range from 16.41 to 23.11 out of 25, so even the lowest-ranked generator received about 66\% of the available points while placing near the bottom of most rankings.



\subsection{Real-POCQi Self-Preference}


Every judge ranked its own answer more favorably than the rest of the panel ranked that same answer. On average a judge placed its own answer 1.22 positions closer to first (95\% CI $[-1.26,-1.18]$), and all eight judge-specific effects remain significant after Holm correction. Correcting for the position in which a candidate was presented leaves this number unchanged, so presentation order does not explain it. In a list of eight an answer can move at most seven positions, so 1.22 is roughly 17\% of the full range; for the strongest judge, GPT-5.6 Sol, it is nearly three positions.

\begin{table}[htbp]
\centering
\small
\caption{Position-adjusted Answer Rank Bias in Real-POCQi. Negative
values indicate that the judge ranked its own answer more favorably.}
\label{tab:real-self}
\begin{tabular}{lrr}
\toprule
Judge & Answer Rank Bias & 95\% CI \\
\midrule
GPT-5.6 Sol & \textbf{-2.94} & $[-3.04,-2.84]$ \\
GPT-5.6 Terra & -2.22 & $[-2.33,-2.12]$ \\
Gemini 3.7 Flash & -1.24 & $[-1.32,-1.16]$ \\
Claude Opus 5 & -0.97 & $[-1.02,-0.92]$ \\
Qwen 3.5 122B & -0.73 & $[-0.85,-0.61]$ \\
Gemini 3.1 Pro & -0.68 & $[-0.79,-0.58]$ \\
Claude Sonnet 5 & -0.65 & $[-0.76,-0.53]$ \\
Qwen 3.8 27B & -0.32 & $[-0.42,-0.23]$ \\
\midrule
Pooled within question & \textbf{-1.22} & $\boldsymbol{[-1.26,-1.18]}$ \\
\bottomrule
\end{tabular}
\end{table}
\FloatBarrier

The effect appears across all judges but varies substantially in size. Judge-specific estimates range from $-0.32$ to $-2.94$, and this ordering does not follow generator quality. The two GPT judges show the largest effects, even though Terra ranks second and Sol fourth in overall performance, while Opus produces the strongest answers overall but shows only a mid-sized effect. This suggests that self-preference is a property of the individual judge rather than simply a result of answer quality. The Qwen judges make this especially clear. As generators, both Qwen models ranked last overall (Table~\ref{tab:real-performance}), with mean ranks of 6.52 and 6.99 out of eight and first-place rates below one percent. A measure based only on wins would therefore show no bias for the Qwen judges. However, each Qwen judge still ranks its own answer higher than the outside judges do.


\begin{table}[htbp]
\centering
\small
\caption{Real-POCQi measures. Answer Score Bias is on the 0--5
mean-rubric scale. Own-Answer Win Rates convert strict listwise rankings to all
own-answer-versus-competitor decisions.}
\label{tab:real-draft-scores}
\begin{tabular}{lrr}
\toprule
Model/judge & Answer Score Bias (95\% CI) & Own-Answer Win Rate \\
\midrule
Qwen 3.5 122B & $+0.64;[+0.60,+0.69]$ & 30.2\% \\
Qwen 3.8 27B & $+0.00;[-0.05,+0.05]$ & 18.8\% \\
Claude Opus 5 & $+0.02;[+0.00,+0.03]$ & \textbf{99.4\%} \\
Claude Sonnet 5 & $-0.09;[-0.12,-0.07]$ & 59.8\% \\
Gemini 3.1 Pro & $+0.55;[+0.52,+0.59]$ & 56.8\% \\
Gemini 3.7 Flash & $+0.67;[+0.65,+0.69]$ & 73.7\% \\
GPT-5.6 Sol & $+0.32;[+0.30,+0.35]$ & 89.5\% \\
GPT-5.6 Terra & $\boldsymbol{-0.75;[-0.79,-0.71]}$ & 93.7\% \\
\bottomrule
\end{tabular}
\end{table}
\FloatBarrier


Answer Score Bias and Own-Answer Win Rate diverge from each other and from Answer Rank Bias (Table~\ref{tab:real-draft-scores}). Terra has the most negative Answer Score Bias of any judge ($-0.75$), yet it ranks its own answer above 93.7\% of competitors. Opus shows almost no score inflation ($+0.02$), yet it places its own answer above 99.4\% of competitors, which is largely consistent with the fact that outside judges also rank Opus first. Qwen 3.5 122B inflates its own scores more than any judge except Gemini Flash but reaches only a 30.2\% Own-Answer Win Rate, because its answers are far enough behind the other candidates that scoring them generously still leaves them near the bottom of the ranking. These measures disagree because they are not measuring the same thing. For instance, Own-Answer Win Rate depends only on where a judge placed its own answer in its own ranking, and never on how other judges viewed that answer. Answer Score Bias compares the same answer across judges, but it does so in raw rubric points, and a generous judge's half-point is not equivalent to a strict judge's. Answer Rank Bias compares the same answer across judges using ranks, which are defined relative to the same candidate list and are therefore unaffected by how a judge uses the scale (this is why it carries the primary analysis).


\subsection{Family-level preference in Real-POCQi}

Judges from the same family ranked their family's two answers 1.27 positions closer to first than judges from other families ranked those same answers (95\% CI $[1.23,1.31]$). This was similar to the exact-model result of 1.22 positions. The family-level pattern therefore did not disappear when we expanded the comparison from one model to both models in its family.

One possible explanation is that models from the same family share writing styles, reasoning patterns, or ideas about what makes an answer good. However, this analysis cannot separate that possibility from exact-model preference, because the same-family average includes the exact model that generated each answer. Distinguishing the two would require more models per family and repeated answers from each model.


\subsection{No-rubric direct-ranking sensitivity}

Because the primary condition supplies five explicit rubric criteria, one possibility is that judges favor answers written in the style those criteria reward. To test this, the direct-ranking condition asks the same eight judges to rank the same blinded responses on overall quality alone, with no criteria provided. It uses a random sample of 100 questions from the Real-POCQi dataset.

Removing the rubric increases the rate of self-preference. The pooled effect is $-1.26$ positions (95\% CI $[-1.36,-1.16]$), which is slightly larger than the $-1.22$ effect observed when the rubric is used. Seven of the eight judges still favor their own answers by an amount that excludes zero. The only exception is Qwen 3.8 27B, with an effect of $-0.16$ positions (95\% CI $[-0.35,0.04]$). It also had the smallest effect in the main condition, and this estimate is based on 100 questions rather than 620. The two conditions still produce very similar rankings overall. For any pair of candidates, the rubric and no-rubric conditions agree on which response is better 91.5\% of the time, and they select the same top answer 87.1\% of the time. However, they produce the exact same full eight-answer ranking only 16.1\% of the time. In other words, removing the rubric does not change which answers judges treat as strong or weak; it only reshuffles answers that are already close in quality.


\subsection{Identity-revealed comparison}

The primary conditions hide generator identity, which means that any self-preference detected there came from the judge recognizing its own output instead of recognizing a model label. The identity-revealed condition tests the opposite setup by naming the generator of every candidate. This condition covers 200 of the 620 questions from the Real-POCQi dataset, judged by the six API judges, giving 1,200 judgments, where all eight candidates are still shown, each labeled with its generator. The blinded comparison uses exactly the same 200 questions and the same six judges, so the only difference between the two conditions is whether names were visible. Presentation-position effects are estimated separately in each condition.


\begin{table}[htbp]
\centering
\small
\caption{Final position-adjusted identity comparison. Positive changes indicate
weaker self-preference when generator names are visible.}
\label{tab:identity}
\begin{tabular}{lrrr}
\toprule
Judge & Revealed ARB & Blinded ARB & Change \\
\midrule
Claude Opus 5 & -1.15 & -1.26 & +0.10 \\
Claude Sonnet 5 & -0.80 & -0.83 & +0.04 \\
Gemini 3.1 Pro & -0.45 & -0.75 & \textbf{+0.30} \\
Gemini 3.7 Flash & -1.18 & -1.32 & +0.14 \\
GPT-5.6 Sol & \textbf{-2.74} & \textbf{-2.83} & +0.08 \\
GPT-5.6 Terra & -2.03 & -2.10 & +0.08 \\
\midrule
Pooled within question & \textbf{-1.39} & \textbf{-1.51} & \textbf{+0.12} \\
\bottomrule
\end{tabular}
\end{table}
\FloatBarrier

The pooled difference between the revealed and blinded conditions is $+0.12$ positions (95\% CI $[+0.06,+0.19]$, nominal $p=0.00019$). Interestingly, showing the generator names makes judges slightly \textit{less} likely to favor their own model, rather than more. After Holm correction across the six judge-specific comparisons, the p-values for Claude Opus and Gemini Pro remain below 0.05, while the other four do not. The overall change is also small compared with the size of the blinded effect. A shift of $+0.12$ relative to the blinded effect of $-1.51$ accounts for roughly eight percent of the bias, and every effect in the identity revealed condition remains negative with an interval that excludes zero. This suggests that self-preference is not mainly driven by seeing the model label. Most of the effect is still present when identities are hidden, which is more consistent with judges responding to features such as a model's style, formatting, or reasoning conventions rather than explicit self-identification.

\subsection{Clinical specialty analysis}

If self preference were dependent on the medical domain, it should appear in some specialties and not others. However, we found that it does not. Every specialty with at least ten questions shows self preference. The effects are also fairly similar across specialties, ranging from $-1.54$ positions in Family Medicine to $-0.91$ in Radiology (which is a difference of about 0.6 positions). By comparison, the differences across judges are more than four times larger, with a range of $-0.32$ to $-2.94$. This means that the size of the effect depends more on which model is judging than on the content of what is being judged.

Specialty accounts for 5.7\% of the variation in self-preference across questions. To see whether this amount was greater than what we might expect by chance, we shuffled the specialty labels while keeping the group sizes the same. The observed value was not unusual based on this comparison ($p=0.19$), and no judge specific specialty test had a $p$-value below 0.05. The smallest specialties are based on too few questions to rule out a real difference, which means that we can only say that this sample shows no clear evidence that self preference depends on specialty.


\subsection{MedSP1000 results}

There were a total of 1,600 judgments at each of the four conversation lengths with the same eight models that acted as generators and judges. At every length, the panel as a whole
ranked its own conversations higher than the outside judges ranked those same
conversations. The result also held after putting each judge's rubric scores
on a common scale (last column of Table~\ref{tab:med-pooled}).

\begin{table}[htbp]
\centering
\small
\caption{Average MedSP1000 self-preference by the number of visible turns.
Negative Answer Rank Bias means that judges placed their own conversation higher
than the other judges did.}
\smallskip
\label{tab:med-pooled}
\begin{tabular}{rrrr}
\toprule
Visible turns & Answer Rank Bias & 95\% CI & Standardized rubric effect \\
\midrule
2 & -0.57 & $[-0.65,-0.49]$ & -0.25 \\
4 & -0.53 & $[-0.60,-0.45]$ & -0.25 \\
6 & -0.62 & $[-0.70,-0.54]$ & -0.30 \\
8 & -0.61 & $[-0.69,-0.53]$ & -0.30 \\
\bottomrule
\end{tabular}
\end{table}
\FloatBarrier


The pooled estimate was smallest at four turns and largest at six turns, but
the differences across lengths were not clear in the overall repeated-measures
test ($F(3,597)=1.69$, $p=0.17$). A paired multivariate test led to the same
conclusion ($p=0.13$). Conversation length therefore did not show a consistent
relationship with self-preference in the eight-model evaluation.


\begin{table}[htbp]
\centering
\small
\caption{Judge-specific MedSP1000 Answer Rank Bias.}
\smallskip
\label{tab:med-judge}
\begin{tabular}{lrrrr}
\toprule
Judge & 2 turns & 4 turns & 6 turns & 8 turns \\
\midrule
Claude Opus 5 & -0.10 & -0.46 & -0.64 & -0.47 \\
Claude Sonnet 5 & -0.04 & -0.35 & -0.39 & -0.37 \\
Gemini 3.1 Pro & \textbf{-1.05} & -0.67 & -0.67 & -0.65 \\
Gemini 3.7 Flash & -0.78 & -0.68 & -0.75 & -0.80 \\
GPT-5.6 Sol & -0.77 & \textbf{-0.73} & \textbf{-0.87} & \textbf{-1.20} \\
GPT-5.6 Terra & -0.87 & -0.72 & -0.86 & -0.90 \\
Qwen 3.5 122B & -0.62 & -0.38 & -0.53 & -0.41 \\
Qwen 3.8 27B & -0.32 & -0.21 & -0.24 & -0.11 \\
\bottomrule
\end{tabular}
\end{table}
\FloatBarrier


All 32 model-by-length estimates were negative. After Holm correction, all
model-specific tests were below 0.05 except Opus and Sonnet at two turns and
Qwen 3.8 27B at eight turns. Models followed different patterns across length,
which helps explain the absence of a consistent pooled trend.

The secondary measures again tell different stories
(Table~\ref{tab:med-draft-scores}).



\begin{table}[htbp]
\centering
\scriptsize
\caption{Two secondary MedSP1000 measures by visible length. Each cell shows
Answer Score Bias on the 0--5 scale, followed by the percentage
of competing conversations it ranked below its own.}
\smallskip
\label{tab:med-draft-scores}
\begin{tabular}{lrrrr}
\toprule
Model/judge & 2 turns & 4 turns & 6 turns & 8 turns \\
\midrule
Claude Opus 5 & $-0.53$ / 51.6\% & $-0.26$ / 60.9\% & $-0.16$ / 72.9\% & $-0.14$ / \textbf{88.4\%} \\
Claude Sonnet 5 & $-0.48$ / 52.0\% & $-0.14$ / \textbf{78.1\%} & $-0.10$ / \textbf{84.6\%} & $-0.10$ / 82.5\% \\
Gemini 3.1 Pro & $\boldsymbol{+0.71}$ / 63.0\% & $+0.49$ / 50.4\% & $+0.39$ / 46.1\% & $+0.25$ / 42.8\% \\
Gemini 3.7 Flash & $+0.68$ / 52.9\% & $\boldsymbol{+0.59}$ / 51.1\% & $\boldsymbol{+0.58}$ / 52.6\% & $\boldsymbol{+0.63}$ / 54.6\% \\
GPT-5.6 Sol & $+0.12$ / 68.4\% & $+0.07$ / 66.9\% & $+0.12$ / 68.7\% & $+0.21$ / 77.7\% \\
GPT-5.6 Terra & $-0.36$ / 59.8\% & $-0.20$ / 56.1\% & $-0.09$ / 58.4\% & $+0.08$ / 58.1\% \\
Qwen 3.5 122B & $+0.52$ / \textbf{72.6\%} & $+0.31$ / 53.9\% & $+0.30$ / 43.9\% & $+0.28$ / 29.9\% \\
Qwen 3.8 27B & $+0.20$ / 37.4\% & $+0.10$ / 34.8\% & $+0.12$ / 34.4\% & $+0.14$ / 27.5\% \\
\bottomrule
\end{tabular}
\end{table}
\FloatBarrier


The Qwen results show why these measures should remain separate. Qwen 3.5 122B
ranked first overall at two turns but fell to eighth at eight turns. Its
Own-Answer Win Rate fell with this change in conversation quality, from 72.6\% to
29.9\%. Nevertheless, its judge continued to rank its own conversation higher
than outside judges ranked that same conversation at every length. Qwen 3.8
27B also had negative Answer Rank Bias estimates at every length, although its
eight-turn interval included zero (95\% CI $[-0.28,+0.06]$). Conversely,
Opus's Own-Answer Win Rate rose from 51.6\% to 88.4\% as its conversations improved,
even though it scored its own work more harshly than outside judges did. The
Own-Answer Win Rate follows both conversation quality and judging behavior;
Answer Rank Bias holds conversation quality fixed.

\subsection{Token-length sensitivity}

Longer responses generally received better evaluations, even after comparing
answers within the same question and judge and accounting for the generator
and answer position (Table~\ref{tab:length}). In Real-POCQi, doubling length
was associated with a 1.31-place improvement in rank and a 0.59-standard-
deviation increase in rubric score. MedSP1000 showed the same directions at
every conversation length. These are associations; they do not show that
adding words would improve an answer.

\begin{table}[htbp]
\centering
\scriptsize
\caption{Changes associated with doubling response length. Negative rank
changes are better. The final column asks whether a model shows more or less
self-preference on its longer-than-usual answers; negative values mean more.
Brackets show 95\% confidence intervals grouped by question.}
\label{tab:length}
\begin{tabular}{lrrr}
\toprule
Condition & Rank change & Rubric score change & ARB change \\
\midrule
Real-POCQi & $-1.31\;[-1.43,-1.20]$ & $+0.59\;[+0.53,+0.64]$ & $\boldsymbol{+0.54\;[+0.45,+0.63]}$ \\
MedSP1000, 2 turns & $-1.66\;[-1.93,-1.39]$ & $+0.68\;[+0.56,+0.80]$ & $-0.30\;[-0.43,-0.17]$ \\
MedSP1000, 4 turns & $-2.83\;[-3.13,-2.53]$ & $+1.15\;[+1.01,+1.29]$ & $-0.33\;[-0.47,-0.19]$ \\
MedSP1000, 6 turns & $\boldsymbol{-3.00\;[-3.30,-2.69]}$ & $\boldsymbol{+1.23\;[+1.09,+1.36]}$ & $-0.29\;[-0.42,-0.16]$ \\
MedSP1000, 8 turns & $-2.70\;[-3.04,-2.36]$ & $+1.15\;[+0.99,+1.30]$ & $-0.24\;[-0.41,-0.07]$ \\
\bottomrule
\end{tabular}
\end{table}
\FloatBarrier

The final column differs across the two tasks. In Real-POCQi, models showed
less self-preference on their longer-than-usual answers. In MedSP1000, they
showed more self-preference on their longer-than-usual conversations. Response
length affects evaluation, but this reversal means it cannot provide one
common explanation for the main finding in both datasets.

\subsection{Clinical validation}

Blinded physician review remains incomplete. The rankings reported here
describe how automated judges compare the answers, not how closely those
rankings agree with clinical experts.




\section{Discussion}
As medical benchmarking and monitoring relies more on LLMs-as-a-judge, it has become increasingly important to understand how LLMs exhibit self-preference. \cite{alvarezarenas2026unreliable} found that clinical LLM judges may favor AI-generated responses over human-written responses, partly because of response length and writing style.
\cite{panickssery2024llm} showed that LLM judges can recognize and favor their own generations on general tasks like news summarization, while \cite{wataoka2025selfpreference} found that judges tend to prefer lower-perplexity responses that feel more familiar to them. Our study focuses on healthcare, where benchmarking results can impact which models are selected for implementation. We examine self-preference across several model families by separating scoring generosity and direct preference for a model’s own response from Answer Rank Bias, which holds the exact answer fixed and changes only the judge. In the single turn setting, we test whether the self preference effect remains when the rubric is removed and when generator identities are shown to the models. We also examine whether self-preference changes as multi-turn patient–physician conversations become longer, from two to eight turns.

\subsection{Single Turn using Real-POCQi}

Our experiments show that self-preference can appear differently depending on how it is measured. Every Real-POCQi judge ranks its own answer more favorably than outside judges rank that same answer. This includes models whose answers rarely rank first. The Qwen judges show why the choice of measure matters. Their answers rank first in fewer than one percent of judgments, so a metric based only on wins would show no bias. However, each Qwen judge still ranks its own answer higher than the rest of the panel does. Own-Answer Win Rate by itself does not separate a judge favoring its own answer from a judge simply having written a better one. This is why we use Answer Rank Bias, which holds the answer fixed. Self-preference persists after accounting for presentation position and scoring generosity. It also persists for seven of the eight judges when the rubric is removed entirely. Because judges never saw who wrote what in the blind condition, the effect is better described as an affinity for a model's own output than as self-recognition. This affinity may come from shared style or preferred kinds of evidence.

In our study removing the rubric did not weaken the self-preference effect. The effect is not due to judges rewarding answers that happen to match a given criteria. Together with prior work showing that self-preference survives fully objective rubrics \citep{pombal2026selfpreference}, this suggests that rubric design is not an effective tool for controlling self-preference.

When the generator names are shown, the judges still favor their own answer relative to the outside panel. In the identity revealed scenario, we would expect self preference to increase. However, we observed that there was a slight reduction in self-preference effect. The reduction results in about eight percent of the bias, which means that hiding model identity (an obvious safeguard against self preference) addresses very little of the effect. Therefore, explicit labels are not necessary for same-model affinity to occur and providing them does not increase it.

The family-level analysis points in the same direction. Judges ranked both answers from their family more favorably than outside judges did, by about the same amount as in the exact-model analysis. This is consistent with shared family-level tendencies. However, the calculation includes the exact-model judge, so it cannot tell us whether the other model in the family independently shows the same preference.

\subsection{Multi Turn using MedSP1000}


In MedSP1000, pooled self-preference appears at every tested conversation length, but the overall test does not show a clear difference across lengths. The largest pooled point estimate occurred at six turns, followed closely by eight turns. Individual judges follow different paths. Most notably, Qwen 3.5 122B continues to move its own conversation upward even as its overall quality rank falls from first at two turns to eighth at eight turns. Answer Rank Bias separates this judging preference from changes in conversation quality. The data do not support a general claim that longer interactions produce more self-preference.

\bigskip

Response length is a separate concern. Longer answers generally receive better evaluations, whether because of added detail, unnecessary verbosity, or both. However, the main comparison holds the exact answer fixed. In addition, models showed less self-preference on their longer Real-POCQi answers but more on their longer MedSP1000 conversations. Length affects evaluation, but it does not provide one explanation for the pattern in both tasks.

We found no clear relationship between self-preference and the clinical specialty of the question. Across specialties with enough questions, the effect varies by about 0.6 rank positions. However, across judges it varies by more than 2.5 rank positions (from $-0.32$ to $-2.94$). What matters is which model is judging, not what the question is about. This is practically important because the two largest effects in our cohort belong to the GPT judges, and OpenAI models are currently the most widely used evaluators in medical AI benchmarking. A benchmark that relies on a single evaluator from that family would favor answers from it, without any indication in the reported scores that this had happened. Medical benchmarking studies should therefore report results by judge and use multiple judge families rather than relying on a single pooled evaluator.


\section{Limitations}


Given the breadth of possible experiments and confounders, there are a few notable limitations to the work. Each model answered each question once and each judging condition was run once, so we cannot tell how much an individual effect would change if the same experiment were repeated. We also did not benchmark the judges against physician evaluations, so in some cases a judge may prefer its own answer because that answer is genuinely better. We try to control for this by including rankings/evaluations by a plethora of models, but follow up human-expert evaluation would strengthen the analysis. We also attempt to deepen understanding of self preference effects in relation to token length, same model-family affinity, and different grading procedures. However, each of these areas was done on subsamples of our initial dataset and would benefit from further analysis. Finally, this paper presents results for the frontier models available at the time of the experiments, and these results may not reflect the highest performing models even a few months from now. To address this, we plan to continuously evaluate self-preference in newer models as they are released, in the form of a public medical self-preference leaderboard.

\section{Conclusion}

Overall, our findings show that self-preference is present in medical LLM judging, but it does not appear in the same way across all models or measures. The experiments show that LLM judges rank answers from their own exact model more favorably than other judges rank those same answers. This pattern remains across two medical settings, after adjusting for presentation position, using rubric scores, and in a no-rubric sensitivity analysis. In multi-turn conversations, pooled self-preference was present at every tested length, without a clear overall difference across lengths. Revealing generator names slightly reduces Answer Rank Bias while still leaving strong same-model affinity. These findings highlight the importance of using multiple judge families, multiple evaluation measures, and fixed-answer analyses when comparing medical LLMs.


\bibliographystyle{unsrtnat}
\bibliography{references}

\end{document}
